\section{Modelo de voz (Omni): la unificación de escuchar y hablar}

\subsection{Arquitectura Thinker-Talker}

Qwen-Omni adopta una arquitectura ``Thinker + Talker'':

\begin{itemize}
    \item \textbf{Thinker}: encargado de la comprensión y el razonamiento
    \item \textbf{Talker}: encargado de la síntesis y la salida de voz
    \item Logra una interacción de voz end-to-end (sin depender de un pipeline de ASR)
\end{itemize}

\begin{knowledgebox}{Voz end-to-end frente a pipeline de ASR}
El enfoque tradicional es un pipeline de tres etapas: ASR → LLM → TTS, pero Qwen-Omni sigue la ruta end-to-end: el modelo comprende directamente la entrada de voz y genera la salida de voz. Esto no solo reduce la latencia, sino que también permite capturar información no textual como la entonación y la emoción. En cuanto al TTS, el timbre de voz se puede personalizar describiendo las características de la voz mediante un prompt.
\end{knowledgebox}

Actualmente, la inteligencia lingüística del modelo Omni presenta una ligera degradación (aproximadamente al nivel de Gemini 2.5 Flash), pero su capacidad de voz está a la par de Gemini 2.5 Pro.

\subsection{Resumen de la sección}

La voz es un canal importante para la interacción entre el agente y los humanos. La arquitectura Thinker-Talker de Qwen-Omni es un paso clave para lograr el objetivo de ``tres entradas y tres salidas'' (la comprensión y la generación de texto, visión y audio).
